\documentclass[12pt,a4paper]{article}

% ------------------------------------------------
% Pacotes Fundamentais
% ------------------------------------------------
\usepackage[utf8]{inputenc}
\usepackage[T1]{fontenc}
\usepackage[brazil]{babel}
\usepackage{graphicx}
\usepackage[table]{xcolor}
\usepackage{array}
\usepackage{tabularx}
\usepackage{geometry}
\usepackage{fancyhdr}
\usepackage{mathptmx} % Times New Roman
\usepackage{enumitem} 
\usepackage{amsfonts}
\usepackage{amsmath}
\usepackage{amssymb}
\usepackage{tikz}
\usetikzlibrary{positioning, shapes.geometric, calc} 

% Cores Institucionais
\definecolor{maincolor}{RGB}{0,100,50}
\definecolor{forestgreen}{RGB}{0,100,50}

% ------------------------------------------------ 
% Configuração de Página e Margens
% ------------------------------------------------
\geometry{
    a4paper,
    left=2cm,
    right=2cm,
    top=2cm,
    bottom=2cm,
    headheight=15pt
}

\pagestyle{fancy}
\fancyhf{}
\renewcommand{\headrulewidth}{0pt}
\setlength{\parindent}{0pt}

% ------------------------------------------------
% Macros Customizadas (Caixa de Resposta e Alinhamento)
% ------------------------------------------------
\newcommand{\resposta}[1]{}

\newcolumntype{L}[1]{>{\raggedright\arraybackslash}p{#1}}
\newcolumntype{C}[1]{>{\centering\arraybackslash}p{#1}}

% ------------------------------------------------
% Configurações de Caixas Decorativas (tcolorbox)
% ------------------------------------------------
\usepackage[most]{tcolorbox}
\newtcolorbox{blockbox}[1]{%
  colback=green!5,
  colframe=forestgreen,
  coltitle=white,
  title=#1,
  fonttitle=\bfseries,
  boxrule=1.5pt,
  arc=3mm,
  left=8pt,
  right=8pt,
  top=8pt,
  bottom=8pt
}

% ------------------------------------------------
% Configurações de Código (Listings Java)
% ------------------------------------------------
\usepackage{listings}
\definecolor{codegray}{rgb}{0.5,0.5,0.5}
\definecolor{codepurple}{rgb}{0.58,0,0.82}
\definecolor{backcolour}{rgb}{0.96,0.96,0.94}

\lstdefinestyle{javastyle}{
    backgroundcolor=\color{backcolour},   
    commentstyle=\color{codegray},
    keywordstyle=\color{forestgreen}\bfseries,
    numberstyle=\tiny\color{codegray},
    stringstyle=\color{codepurple},
    basicstyle=\ttfamily\footnotesize,
    breakatwhitespace=false,         
    breaklines=true,                 
    captionpos=b,                    
    keepspaces=true,                 
    numbers=none,                    
    showspaces=false,                
    showstringspaces=false,
    showtabs=false,                  
    tabsize=4,
    language=Java,
    literate={á}{{\'a}}1 {ã}{{\~a}}1 {â}{{\^a}}1 {à}{{\`a}}1 {é}{{\'e}}1 {ê}{{\^e}}1 {í}{{\'i}}1 {ó}{{\'o}}1 {õ}{{\~o}}1 {ô}{{\^o}}1 {ú}{{\'u}}1 {ç}{{\c{c}}}1 {Á}{{\'A}}1 {Ã}{{\~A}}1 {Â}{{\^A}}1 {À}{{\`A}}1 {É}{{\'E}}1 {Ê}{{\^E}}1 {Í}{{\'I}}1 {Ó}{{\'O}}1 {Õ}{{\~O}}1 {Ô}{{\^O}}1 {Ú}{{\'U}}1 {Ç}{{\c{C}}}1
}
\lstset{style=javastyle}

% ================================================
\begin{document}
% ================================================

% ------------------------------------------------
% Cabeçalho Institucional Oficial da UniRV
% ------------------------------------------------
\begin{tabularx}{\textwidth}{@{} L{4.5cm} X @{}}
    \raisebox{-0.5cm}{\includegraphics[width=4.5cm]{Universidade_de_Rio_Verde.jpg}} &
    \hfill
    \begin{tabular}{|>{\centering\arraybackslash}m{11.6cm}|}
        \hline
        \cellcolor{forestgreen}\color{white}\rule{0pt}{1.2cm}\textbf{ROTEIRO PRÁTICO DA DISCIPLINA} \tabularnewline
        \hline
    \end{tabular} 
\end{tabularx}
\vspace{0.3cm}

\begin{tabularx}{\textwidth}{|X|L{2.5cm}|}
\hline
\textbf{Curso} & \textbf{Campus} \\
\textit{Engenharia de Software} & \textit{Rio Verde} \\
\hline
\multicolumn{2}{|l|}{\textbf{Disciplina}} \\  
\multicolumn{2}{|l|}{ \textit{Programação Orientada a Objetos}} \\
\hline
\textbf{Professor} & \textbf{Turma} \\
\textit{Sergio Souza Novak} & \textit{A / B} \\
\hline
\end{tabularx}

\vspace{0.5cm}

\section*{Tratamento de Exceções em Java}

Um ótimo jeito de aprender exceções em Java é modelar situações reais de uma conta corrente, porque vários erros fazem sentido no domínio do problema:
\begin{itemize}
    \item saldo insuficiente
    \item valor inválido
    \item conta inexistente
    \item saque acima do limite
    \item transferência para a mesma conta
\end{itemize}

Isso ajuda a entender:
\begin{itemize}
    \item \texttt{throw}
    \item \texttt{throws}
    \item exceções personalizadas
    \item \texttt{try/catch}
    \item checked vs unchecked exceptions
\end{itemize}

\section*{Roteiro Prático}

\subsection*{1. Criar a classe da conta}
\begin{lstlisting}
public class ContaCorrente {

    private String titular;
    private double saldo;

    public ContaCorrente(String titular, double saldoInicial) {
        this.titular = titular;
        this.saldo = saldoInicial;
    }

    public double getSaldo() {
        return saldo;
    }
}
\end{lstlisting}

\subsection*{2. Criar uma exceção personalizada}
Exemplo: saldo insuficiente
\begin{lstlisting}
public class SaldoInsuficienteException extends Exception {

    public SaldoInsuficienteException(String mensagem) {
        super(mensagem);
    }
}
\end{lstlisting}

Aqui você aprende:
\begin{itemize}
    \item herança
    \item \texttt{extends Exception}
    \item construtor de exceção
\end{itemize}

\subsection*{3. Lançar exceção com \texttt{throw}}

Agora vamos criar o saque.

\begin{lstlisting}
public void sacar(double valor)
        throws SaldoInsuficienteException {

    if (valor > saldo) {
        throw new SaldoInsuficienteException(
                "Saldo insuficiente para saque."
        );
    }

    saldo -= valor;
}
\end{lstlisting}

\subsection*{4. Tratar a exceção com \texttt{try/catch}}
\begin{lstlisting}
public class Main {

    public static void main(String[] args) {

        ContaCorrente conta =
                new ContaCorrente("Ana", 500);

        try {

            conta.sacar(1000);

        } catch (SaldoInsuficienteException e) {

            System.out.println("Erro: " + e.getMessage());

        }

        System.out.println("Programa continua...");
    }
}
\end{lstlisting}

\textbf{Fluxo do que acontece}
\begin{verbatim}
main()
  ↓
sacar()
  ↓
throw new SaldoInsuficienteException()
  ↓
catch()
\end{verbatim}

\subsection*{5. Criar mais exceções}
\textbf{Valor inválido}
\begin{lstlisting}
public class ValorInvalidoException
        extends RuntimeException {

    public ValorInvalidoException(String mensagem) {
        super(mensagem);
    }
}
\end{lstlisting}

\subsection*{6. Trabalhar checked vs unchecked}
\textbf{Checked Exception}
\begin{itemize}
    \item \texttt{extends Exception}
    \item Java obriga: \texttt{try/catch} ou \texttt{throws}
\end{itemize}

\textbf{Unchecked Exception}
\begin{itemize}
    \item \texttt{extends RuntimeException}
    \item Java NÃO obriga tratamento.
\end{itemize}

\textbf{Exemplo prático}
\begin{lstlisting}
public void depositar(double valor) {

    if (valor <= 0) {
        throw new ValorInvalidoException(
                "Valor inválido."
        );
    }

    saldo += valor;
}
\end{lstlisting}

\subsection*{7. Transferência entre contas}
Excelente exercício para praticar propagação de exceções.
\begin{lstlisting}
public void transferir(
        ContaCorrente destino,
        double valor)
        throws SaldoInsuficienteException {

    this.sacar(valor);
 
    destino.depositar(valor);
}
\end{lstlisting}

\subsection*{8. Exceções Nativas Comuns no Java}
Além das exceções personalizadas, a linguagem Java possui várias exceções nativas que ocorrem frequentemente no dia a dia. Você não precisa criá-las, mas é fundamental saber como evitá-las ou tratá-las.

\begin{itemize}
    \item \textbf{\texttt{NullPointerException}}: Ocorre quando você tenta acessar um método ou atributo de um objeto que está com referência nula.
\begin{lstlisting}
ContaCorrente conta = null;
conta.getSaldo(); // Lança NullPointerException
\end{lstlisting}

    \item \textbf{\texttt{IllegalArgumentException}}: Usada quando um método recebe um argumento inadequado. Em vez de criar a \texttt{ValorInvalidoException}, muitas vezes podemos usar esta exceção nativa para validações:
\begin{lstlisting}
public void depositar(double valor) {
    if (valor <= 0) {
        throw new IllegalArgumentException("Valor inválido.");
    }
    saldo += valor;
}
\end{lstlisting}

    \item \textbf{\texttt{ArithmeticException}}: Erros de matemática, como divisão por zero.
\begin{lstlisting}
int resultado = 10 / 0; // Lança ArithmeticException
\end{lstlisting}

    \item \textbf{\texttt{ArrayIndexOutOfBoundsException}}: Tentativa de acessar uma posição que não existe em um array.
\begin{lstlisting}
int[] numeros = {1, 2, 3};
System.out.println(numeros[5]); // Lança ArrayIndexOutOfBoundsException
\end{lstlisting}
\end{itemize}

\subsection*{9. Exercícios recomendados}
\begin{itemize}
    \item \textbf{Básico:} saque simples, depósito, saldo insuficiente.
    \item \textbf{Intermediário:} limite diário, conta bloqueada, valor negativo.
    \item \textbf{Avançado:} transferência, logs de erro, hierarquia de exceções, múltiplos catches, \texttt{finally}.
\end{itemize}

\textbf{Exemplo com \texttt{finally}}
\begin{lstlisting}
try {
    conta.sacar(1000);
} catch (SaldoInsuficienteException e) {
    System.out.println(e.getMessage());
} finally {
    System.out.println("Operação encerrada.");
}
\end{lstlisting}

\subsection*{Estrutura ideal do projeto}
\begin{verbatim}
src/
 |-- Main.java
 |-- ContaCorrente.java
 \-- exceptions/
      |-- SaldoInsuficienteException.java
      \-- ValorInvalidoException.java
\end{verbatim}

\subsection*{Conceitos que você aprende nesse exemplo}
\begin{table}[h]
\centering
\begin{tabular}{|l|l|}
\hline
\textbf{Conceito} & \textbf{Onde aparece} \\
\hline
\texttt{throw} & lançar exceção \\
\texttt{throws} & propagar exceção \\
\texttt{try/catch} & tratamento \\
\texttt{finally} & encerramento \\
Herança & exceções personalizadas \\
Polimorfismo & RuntimeException vs Exception \\
Encapsulamento & saldo privado \\
Regras de negócio & validações bancárias \\
\hline
\end{tabular}
\end{table}

\end{document}
